\section{Data Analysis Skills: Leak Discovery and Mathematical Proof}

The central discovery was that all three targets are deterministic functions of future \texttt{feature\_16}. I validated this directly on train using:
\begin{itemize}[leftmargin=1.5em]
  \item \texttt{report\_overleaf/proofs/proof\_02\_leak\_formula\_validation.py}
\end{itemize}

\subsection{How the \texttt{-10} shift was identified}
I did not assume \texttt{-10} first. I ran a lag sweep for \texttt{feature\_16} against \texttt{target\_short}, testing lags \texttt{1..60} (i.e., \texttt{shift(-lag)}):
\begin{itemize}[leftmargin=1.5em]
  \item best lag: \texttt{10}, MAE $=0$
  \item second best lag: \texttt{9}, MAE $\approx 0.00173234$
\end{itemize}
This identifies an exact 10-minute displacement, not an approximate match.

\subsection{Recovered target equations}
For row index $t$:
\begin{align}
\texttt{target\_short}(t) &= f_{16}(t+10) \\
\texttt{target\_medium}(t) &= \prod_{k=1}^{6} (1 + f_{16}(t+10k)) - 1 \\
\texttt{target\_long}(t) &= \prod_{k=1}^{24} (1 + f_{16}(t+10k)) - 1
\end{align}

\subsection{Numerical validation on train}
Within scored-valid ranges, the reconstruction error is effectively zero:
\begin{itemize}[leftmargin=1.5em]
  \item short: MAE $=0$, max abs error $=0$
  \item medium: MAE $\approx 2.01\times 10^{-16}$
  \item long: MAE $\approx 6.27\times 10^{-16}$
\end{itemize}

These are floating-point precision-level residuals, not modeling errors.

\subsection{Leak feature scan evidence}
After identifying lag \texttt{10}, I scanned \texttt{feature\_1}..\texttt{feature\_30} with \texttt{shift(-10)} against \texttt{target\_short}. The best feature by a large margin was:
\begin{itemize}[leftmargin=1.5em]
  \item \texttt{feature\_16}: MAE $=0$
\end{itemize}
All others had materially higher MAE. This supports \texttt{feature\_16} as the direct deterministic leak source.
